% !TeX root = ../main.tex
% =========================================================================
% 5. MARCO HIPOTÉTICO O PREGUNTAS DE INVESTIGACIÓN
% Pregunta general y tres preguntas específicas (P1 diagnóstico/catálogos, P2 arquitectura, P3 evaluación medible)
% =========================================================================
\section{Marco hipotético o preguntas de investigación}

En concordancia con la matriz metodológica institucional UPS P69 y los objetivos específicos planteados, se formulan las preguntas de sistematización del problema y las correspondientes hipótesis de ingeniería que guiarán la investigación experimental.

\subsection{Preguntas de investigación}

\noindent\textbf{Pregunta General (PG).} ¿En qué medida un agente inteligente para la atención comercial y cotización de materiales eléctricos en Guayaquil, basado en modelos generativos y recuperación relacional determinista, mejora la exactitud de precios, la latencia de respuesta y la eficiencia de recursos computacionales en el proceso de ventas de la distribuidora Power Electrimk?

\vspace{0.3cm}
\noindent\textbf{Pregunta Específica 1 (P1).} ¿Cuáles son los requerimientos comerciales y qué estructura de normalización de datos en catálogos técnicos multimarca se requiere para soportar consultas deterministas de inventario y cotización en la distribuidora Power Electrimk?

\vspace{0.3cm}
\noindent\textbf{Pregunta Específica 2 (P2).} ¿Cómo diseñar e implementar una arquitectura de microservicio ligera en .NET 8 que orqueste la mensajería asíncrona de WhatsApp Cloud API, la inferencia generativa de Gemini y la recuperación local en SQLite preservando la memoria conversacional comercial bajo un consumo contenido de recursos?

\vspace{0.3cm}
\noindent\textbf{Pregunta Específica 3 (P3).} ¿En qué medida la incorporación del mecanismo de recuperación determinista local mitiga la tasa de alucinación, optimiza la exactitud fáctica de cotizaciones y reduce la latencia en percentil P95 frente al modelo generativo operando de manera aislada sin recuperación contextual?

\subsection{Hipótesis de investigación}

\noindent\textbf{Hipótesis General (HG).} La implementación de un agente inteligente para la atención comercial y cotización de materiales eléctricos en Guayaquil con recuperación fáctica local sobre catálogos estructurados en SQLite permite optimizar el embudo de ventas técnicas en Power Electrimk, alcanzando una exactitud de precios y referencias técnicas superior al 95\%, una tasa de alucinación fáctica nula (0\%), y una latencia en percentil P95 inferior a 4.0 segundos, con un perfil de consumo de memoria RAM inferior a 250 MB en infraestructura virtual de entrada en la nube.

\vspace{0.3cm}
\noindent\textbf{Hipótesis Específica 1 (H1).} El levantamiento y la normalización relacional de los atributos esenciales del catálogo multimarca estructurado en SQLite permiten resolver consultas de inventario comercial con tiempos de ejecución en base de datos ($T_{\text{sql}}$) inferiores a 100 milisegundos con 100\% de integridad en campos críticos (código de fábrica, precio y marca).

\vspace{0.3cm}
\noindent\textbf{Hipótesis Específica 2 (H2).} Una arquitectura asíncrona desacoplada en .NET 8 con persistencia de sesiones comerciales permite procesar las ráfagas de webhooks entrantes de WhatsApp Cloud API sin bloqueos de concurrencia, manteniendo la utilización de memoria del proceso por debajo de 250 megabytes.

\vspace{0.3cm}
\noindent\textbf{Hipótesis Específica 3 (H3).} El agente comercial con recuperación fáctica determinista exhibirá una tasa de alucinación significativamente menor ($H_{\text{rate}} = 0\%$, $p < 0{,}05$) y una exactitud de cotización ($Acc_{\text{cot}} \ge 95\%$) estadísticamente superior en comparación con el modelo generativo operando en modalidad directa sin recuperación contextual.
